\documentclass[10pt,a4paper]{amsart}
\usepackage[margin=19mm]{geometry}
\usepackage{amsmath,amssymb,mathtools}
\usepackage[scr=rsfs,cal=euler]{mathalfa}
\usepackage{xfrac}
\usepackage[unicode,hidelinks,bookmarks=false]{hyperref}
\newcommand{\ie}{i.e.\ }
\newcommand{\cf}{cf.\ }
\newcommand{\resp}{respectively\ }
\newcommand{\Subsection}{\subsection}
\providecommand{\nobreakdash}{\nobreak\hbox{-}}
\setlength{\emergencystretch}{2em}
\multlinegap0pt
\usepackage{amssymb}
\usepackage{mathtools}
\usepackage[shortlabels]{enumitem}
\usepackage{xcolor}
\newtheorem{lemma}{Lemma}
\newtheorem{theorem}{Theorem}
\newcommand{\F}{\mathbb F}
\newcommand{\Tau}{\mathcal T}
\newcommand{\prk}{\operatorname{prk}}
\newcommand{\srk}{\operatorname{srk}}
\title{Slice and Partition Rank Criteria for Polynomial Zero-Avoidance}
\author{Simone Costa, Stefano Della Fiore, Mattia Fontana}
\date{Source: arXiv:2607.29490v1 (CC BY 4.0)}
\begin{document}
\maketitle

\noindent\emph{Source and licence.} Slice and Partition Rank Criteria for Polynomial Zero-Avoidance. Version arXiv:2607.29490v1 (CC BY 4.0), distributed by arXiv under the Creative Commons Attribution 4.0 International licence, http://creativecommons.org/licenses/by/4.0/, as recorded in the arXiv licence metadata for this exact version and shown on the abstract page https://arxiv.org/abs/2607.29490v1. Full licence notice: \texttt{licenses/arxiv-2607.29490-CC-BY-4.0.txt}.\par\smallskip

\noindent\textit{Editorial scope.} Counted unit: the article's theorem thm:distinct-reduction (source0.tex lines 1034--1054). The article's complete original proof of it is delivered with it (source0.tex lines 1056--1109, one \texttt{proof} environment of 54 lines). No mathematics is written, repaired or re-derived: the statement and the proof are the article's own lines, in the article's own notation, and no retained line is substituted or re-flowed.  Retained uncounted as the source-local context the proof needs: lines 903--919; lines 971--978.  Nothing was omitted from any retained span, so the package carries no omission record.  No retained line contains a citation, so the wrapper carries no bibliography and no bibliography entry.  No retained line contains a figure environment, an image-inclusion command or drawing code, so no image is delivered and the package declares no asset dependency.  Editorial formatting only, in addition: an AMS \texttt{amsart} preamble that restates the article's own declarations and macros used by the retained lines, namely the environments \texttt{lemma}, \texttt{theorem} on separate counters, together with the standard packages those lines need.  The retained environments print in this document as Lemma 1, Theorem 1 in order of appearance, whereas the article prints the same items with the numbering of its own full document.  The complete native source of the article as distributed in the arXiv e-print for this exact version is bundled unchanged as \texttt{sources/206492/source0.tex}.
\begin{lemma}[Naslund's distinctness indicator]\label{lem:naslund-indicator}
For every $\mathbf x=(x_1,\ldots,x_t)$,
\[
 H_t(\mathbf x)=
 \begin{cases}
 1,&x_1,\ldots,x_t\text{ are pairwise distinct},\\
 (-1)^t(t-1)!,&x_1=\cdots=x_t,\\
 0,&\text{otherwise}.
 \end{cases}
\]
Moreover,
\[
 H_t=\sum_{\pi\in\Pi_t\setminus\{\widehat1\}}c_\pi\Delta_\pi,
\qquad
 c_\pi=(-1)^{t-|\pi|}\prod_{B\in\pi}(|B|-1)!.
\]
\end{lemma}

\begin{lemma}[Lifting contractions]\label{lem:lifting}
If $\pi$ has at least two blocks, then
\[
 \prk(\Delta_\pi\Tau_{F,n})
 \le\prk(\Tau_{F,\pi,n})
 \le\srk(\Tau_{F,\pi,n}).
\]
\end{lemma}

\begin{theorem}[Distinct-variable reduction via contractions]\label{thm:distinct-reduction}
Let $q$ have characteristic $p$, assume $2\le t\le p$, and let $F\in\F_q[x_1,\ldots,x_t]$ satisfy
\[
 F(a,\ldots,a)=0\qquad(a\in\F_q).
\]
If $A\subseteq\F_q^n$ contains no $t$ pairwise distinct elements satisfying $F=0$ coordinatewise, then
\[
 |A|
 \le
 \sum_{\pi\in\Pi_t\setminus\{\widehat1\}}
 \srk(\Tau_{F,\pi,n}).
\]
Consequently, if a sequence $S$ contains no length-$t$ subsequence on which $F$ vanishes, then
\[
 |S|
 \le
 (t-1)
 \sum_{\pi\in\Pi_t\setminus\{\widehat1\}}
 \srk(\Tau_{F,\pi,n}).
\]
\end{theorem}

\begin{proof}
Consider
\[
 J=H_t\Tau_{F,n}.
\]
On $A^t$, a tuple of pairwise distinct elements has
$\Tau_{F,n}=0$ by hypothesis; a partial diagonal has $H_t=0$;
and a constant tuple has
\[
 J(x,\ldots,x)=(-1)^t(t-1)!.
\]
Since $t\le p$, this scalar is nonzero in $\F_q$. Hence
$J|_{A^t}$ is diagonal with nonzero diagonal. By Naslund's
partition-rank diagonal lemma,
\[
 |A|\le\prk(J).
\]

By Lemma~\ref{lem:naslund-indicator}, with the integer coefficients
$c_\pi$ viewed in $\F_q$,
\[
 J
 =H_t\Tau_{F,n}
 =\sum_{\pi\in\Pi_t\setminus\{\widehat1\}}
 c_\pi\Delta_\pi\Tau_{F,n}.
\]
{Subadditivity of partition rank gives
\[
 \prk(J)
 \le
 \sum_{\pi\in\Pi_t\setminus\{\widehat1\}}
 \prk(c_\pi\Delta_\pi\Tau_{F,n})
 \le
 \sum_{\pi\in\Pi_t\setminus\{\widehat1\}}
 \prk(\Delta_\pi\Tau_{F,n}),
\]
because multiplication by a nonzero scalar does not change partition
rank, while if $c_\pi=0$ in $\F_q$ the corresponding term on the
left is zero. Applying Lemma~\ref{lem:lifting}, we obtain
\[
 |A|
 \le\prk(J)
 \le\sum_{\pi\ne\widehat1}\prk(\Delta_\pi\Tau_{F,n})
 \le\sum_{\pi\ne\widehat1}\srk(\Tau_{F,\pi,n}).
\]
Let \(A\) be the set of distinct values appearing in \(S\).
Since no value can occur \(t\) times, we have
\[
 |A|\ge \frac{|S|}{t-1}.
\]
If \(A\) contained \(t\) distinct elements satisfying \(F=0\),
choosing one occurrence of each in \(S\) would give a forbidden
subsequence. Apply the set bound to \(A\).}
\end{proof}

\end{document}
